\documentclass{osglecture}
\input{seriesexample-exercises}

\begin{document}
\lecture[\lende{Exercises}{Übungen}]
  {\lende{Exercises on the shared source}{Übungen zur gemeinsamen Quelle}}
  {exercises}

\section{\lende{Exercises}{Aufgaben}}
\IfLectureModeTF{solutions}
  {\lende{This version contains the solutions.}%
    {Diese Fassung enthält die Lösungen.}}
  {\lende{The solutions are published separately.}%
    {Die Lösungen werden gesondert veröffentlicht.}}

\begin{exercise}
  \lende{Name the two documents that the target \emph{script} and the target
    \emph{slides} produce from the unit introduced in
    Section}{Nennen Sie die beiden Dokumente, die die Targets \emph{script}
    und \emph{slides} aus der in Abschnitt}~\olref[introduction]{sec:shared-source}%
  \lende{.}{\ eingeführten Unit erzeugen.}
\end{exercise}
\begin{solution}
  \lende{A long-form report and a presentation, both from the same
    source.}{Eine Langform und eine Präsentation, beide aus derselben
    Quelle.}
\end{solution}

\begin{exercise}
  \lende{Why does this unit build neither \emph{slides} nor a presentation
    at all?}{Warum baut diese Unit weder \emph{slides} noch sonst eine
    Präsentation?}
\end{exercise}
\begin{solution}
  \lende{Its directory carries the scope code \texttt{a}; OLLM builds it
    only for targets whose \texttt{unit\_scopes} contain \texttt{a}, here
    \emph{script} and \emph{solutions}.}{Ihr Verzeichnis trägt den
    Scope-Code \texttt{a}; OLLM baut sie nur für Targets, deren
    \texttt{unit\_scopes} \texttt{a} enthalten, hier \emph{script} und
    \emph{solutions}.}
\end{solution}
\end{document}
